\section{\revadd{Related Work}}
\subsection{Object Detection and Instance Segmentation}
Accurate object perception is the foundation of robotic manipulation. Traditional 2D detectors such as Faster R-CNN \cite{ren2015faster} and 3D detectors such as PointNet++ \cite{qi2017pointnet++} and PV-RCNN \cite{shi2020pv} primarily output bounding boxes. Although bounding boxes are computationally efficient, they provide only coarse geometric support and can introduce significant extrapolation errors when baggage items overlap or touch each other. Instance segmentation methods, including Mask R-CNN \cite{he2017mask}, YOLACT \cite{bolya2019yolact}, and YOLO-based segmentation models \cite{Jocher_Ultralytics_YOLO26_Unified_2026}, provide pixel-level object separation and are therefore more suitable for downstream point-cloud extraction. Foundation models such as Segment Anything Model (SAM) \cite{kirillov2023segment} show strong generalization, but their computational cost is usually too high for embedded real-time robotic control. In this paper, a lightweight YOLOv10-based architecture is modified for baggage segmentation so that real-time inference can be maintained while the tag-interference problem is addressed.

\subsection{Geometric Modeling}

After obtaining the 2D target mask, robotic grasping requires accurate 3D geometric modeling to estimate the object pose and dimensions. Classical pipelines usually rely on the Point Cloud Library (PCL) \cite{rusu20113d}, which provides standard operations such as VoxelGrid downsampling, statistical outlier removal, and RANSAC-based plane segmentation \cite{fischler1981random}. The cleaned point cloud is then commonly processed by Principal Component Analysis (PCA) to estimate the principal axes and construct an oriented bounding box (OBB). However, baggage-handling scenarios introduce several adverse sensing conditions. Reflective stickers, luggage tags, flexible straps, and conveyor motion may produce holes, ghost points, and local geometric distortions in depth observations from ToF or structured-light cameras. Therefore, depth completion and point cloud denoising have become important techniques for robust 3D geometric modeling.

Depth completion methods aim to recover missing geometry from incomplete depth measurements. For example, HDCNet \cite{xie2025hdcnet} adopts a hybrid architecture combining CNN, Transformer, and Mamba modules to infer the geometry of transparent or reflective objects from RGB and depth cues. Such methods can compensate for missing depths caused by specular reflection and reduce hole-induced pose errors. Nevertheless, they typically require large-scale training data and introduce considerable computational overhead. In baggage handling, reflective regions often appear only locally, such as on stickers or tags, while the main suitcase body is usually close to a regular cuboid. Therefore, a lightweight heuristic treatment is more suitable for real-time deployment.

Point cloud denoising has also been widely studied. Traditional filters, such as bilateral filtering \cite{digne2017bilateral, si2022lbf}, can suppress isolated outliers but are less effective for flexible accessories whose geometry does not follow smooth rigid surfaces. Graph-based methods further exploit local geometric similarity, such as feature graph learning with Graph Laplacian regularization \cite{hu2020feature}. However, these methods still rely heavily on local surface consistency assumptions, which may fail for thin tags and deformable straps. Recent learning-based denoising methods, including Total Denoising \cite{hermosilla2019total, rakotosaona2020pointcleannet}, and EC-Net \cite{yu2018ec}, improve general denoising performance but are not specifically designed for structured baggage accessories and often require additional computation.

Dynamic point cloud denoising further considers temporal inconsistency caused by object or sensor motion. Hu et al. proposed to use spatio-temporal convolution to aggregate gradient fields across frames \cite{hu2025dynamic}, while Hu et al. exploited intra-frame self-similarity and inter-frame consistency through spatio-temporal graph learning \cite{hu2021dynamic}. These approaches are effective for general dynamic point clouds, but they are usually designed for sensor motion or high-speed targets. In our setting, the camera is fixed and the conveyor speed is controllable. Therefore, instead of employing a heavy dynamic denoising network, this work introduces a task-specific Pixel-to-3D Cross Consistency Conditional Filtering (PX3C) module to remove flexible tags and straps, and a Consistency-Aware OBB Stabilization (CAOS) module to suppress frame-to-frame pose jitter. This design preserves the efficiency of classical PCL-based processing while improving robustness in practical baggage-grasping scenarios.

\subsection{Dynamic Grasping}

Grasp planning can be broadly categorized into static and dynamic approaches. Static grasp synthesis assumes that the target object remains stationary during perception and execution, and has been extensively studied in data-driven grasping literature \cite{bohg2013data}. However, this assumption is not suitable for conveyor-based baggage handling, where objects continuously move, collide, and may undergo pose changes during robot execution. Therefore, robust state estimation and time-aware motion planning are essential for dynamic grasping.
\revadd{Classical dynamic grasping commonly uses visual servoing to update robot motion from visual feedback \cite{chaumette2008visual, wang2023visual, he2023image}. Image-based visual servoing (IBVS) regulates image-plane errors without necessarily reconstructing 3D geometry, but does not directly enforce the metric alignment required for precision placement. Position-based visual servoing (PBVS) regulates 3D pose errors and supports explicit trajectory planning \cite{lippiello2007position}, but depends on reliable pose estimation \cite{zhu2024position}. In conveyor-based baggage handling, occlusion, flexible attachments, and continuous motion degrade visual tracking and pose estimation, while downstream placement requires accurate dimensions and alignment. These conditions motivate geometric refinement and predictive state estimation to support robust visual feedback.}

\revadd{Other methods combine target-state estimation with explicit grasp planning.} For moving target tracking, recursive Bayesian filtering provides a standard and efficient framework. Learning-based trajectory prediction methods, such as Social LSTM \cite{alahi2016social}, can model complex nonlinear motions, but they usually require large-scale training data and are unnecessary when the target motion is physically constrained. In conveyor scenarios, baggage motion is mainly governed by belt velocity with moderate sensing noise. Thus, a discrete Kalman filter with a constant-acceleration motion model \cite{li2003survey} offers a favorable trade-off between robustness and computational efficiency. For dynamic interception, Cowley et al. proposed a perception-and-planning framework for picking moving objects \cite{cowley2013perception}. \revadd{Chen et al. further combine global-to-local visual detection with static-to-dynamic planning for moving objects in dense conveyor clutter \cite{chen2025adaptive}.} \revadd{Hybrid closed-loop methods also incorporate learned components into explicit online planning. VFAS-Grasp refines grasp candidates using visual feedback and uncertainty-aware adaptive sampling \cite{piacenza2024vfas}, whereas Grasp-MPC combines a learned visual value function with model predictive control for reactive grasp execution \cite{yamada2026graspmpc}.} \revadd{For conveyor interception,} we adopt a two-stage hierarchical strategy that decouples fast global interception in Cartesian space from local synchronization in relative velocity space, thereby balancing execution speed and grasping precision. \revadd{Beyond adaptive observation and grasp planning, our framework addresses attachment-induced geometric errors and precision placement through OBB refinement and state regulation.}

 \revadd{More recently, end-to-end perception-to-action methods have learned mappings from visual observations to robot actions. DextrAH-G learns a depth-to-action dexterous arm--hand policy with geometric fabrics for continuous visual feedback.} \revadd{Language-conditioned extensions include} RT-2 \cite{zitkovich2023rt}, the Physical Intelligence $\pi$ series \cite{black2024pi_0, intelligence2025vision}, and OpenVLA \cite{kim2024openvla}. These models demonstrate strong potential in semantic reasoning and cross-task generalization. \revadd{Large model size and inference latency remain deployment concerns for these generalist policies. Real-time chunking addresses latency through asynchronous execution of diffusion- or flow-based action policies.} \revadd{Airport baggage handling additionally requires accurate object dimensions and poses for precision placement. Compared with purely end-to-end visuomotor policies, an explicit perception-and-planning structure supports independent geometric error inspection, motion constraints, and structured replanning during exception recovery. We therefore combine instance segmentation, geometric refinement, and predictive planning to retain interpretability and controllability within embedded real-time constraints.}

\subsection{\revadd{Active State Regulation}}

\revadd{Occlusion and clutter can reduce confidence in object pose estimates, while unfavorable object poses can make grasp execution difficult and increase control uncertainty. Active state regulation addresses these conditions by adjusting object poses toward configurations that are easier to observe and grasp. A common approach is robotic rearrangement, which reconfigures cluttered scenes to facilitate target manipulation \cite{batra2020rearrangement}. Such actions can also improve state estimation by revealing occluded geometry, as studied in interactive perception \cite{bohg2017interactive}.}
\revadd{Wang et al. jointly learn pushing and grasping policies through self-supervision to improve target accessibility in dense clutter \cite{wang2024joint}. Hu et al. exploit equivariant models and grasp score optimization to coordinate pushing and grasping for target retrieval \cite{hu2025pushgrasp}. Lee and Kim plan sequences of prehensile and non-prehensile rearrangement actions to retrieve targets in confined workspaces while minimizing the number of relocations \cite{lee2025rearrangement}.}

DexSinGrasp \cite{xu2025dexsingrasp} and OPG-Policy \cite{ding2024opg} use dexterous hands or end-effectors to separate the target from surrounding clutter during grasping. Murray et al. further proposed an interactive visual failure prediction strategy, where probing actions such as lifting or shaking are used before grasp execution \cite{murray2024learning}.

\revadd{GraspView selects informative viewpoints and aligns multi-view RGB reconstructions for cluttered grasping \cite{wang2025graspview}. Here, the camera remains fixed and active state regulation changes the baggage configuration through suction-assisted relocation, followed by renewed perception and planning.}

\revadd{These studies demonstrate the utility of active rearrangement for grasping in clutter. However, pushing and multistep rearrangement can introduce additional disturbances and delay when large baggage moves continuously through a narrow conveyor workspace. We therefore use uncertainty-triggered suction-assisted relocation as a low-complexity state reset, rather than a general rearrangement policy. The relocation aims to restore observability, after which the system acquires new observations, re-estimates the target geometry and motion, and replans the grasp.}
